\documentclass[12pt]{article}
\usepackage[utf8]{inputenc}
\usepackage[T1]{fontenc}
\usepackage[brazil]{babel}
\usepackage[a4paper,margin=2.5cm]{geometry}
\usepackage{enumitem}
\setlist[itemize]{noitemsep,topsep=4pt}
\title{Sprint 4 - S4\_12\\Quick Sort}
\author{Matheus Moreli Durão Silva\\ Matrícula 12621CBS217\\ Cibersegurança - Bacharelado - Noturno}
\date{2026}
\begin{document}
\maketitle

\section*{Entregável e checklist}
Este documento registra o card S4\_12, o conteúdo estudado, a aplicação no Inventário de Ativos e a forma de conferência. O arquivo complementar é \texttt{quick\_sort.py}. A entrega mantém o escopo introdutório: código claro, funções pequenas e nenhuma dependência que não seja necessária.

\begin{itemize}
\item Conceito: Quick Sort.
\item Evidência: documento e, quando aplicável, exemplo executável na mesma pasta.
\item Ligação com o trabalho: referência direta aos arquivos do repositório principal.
\item Conferência: leitura do exemplo, execução dos testes e comparação com a explicação abaixo.
\end{itemize}

\section*{O que foi estudado e aplicado}
Quick Sort escolhe um pivô, separa menores, iguais e maiores e ordena as partes. O exemplo evita alterar a lista original para facilitar a conferência.

\section*{Lógica da decisão}
O card compara estratégias de ordenação. Não há motivo para substituir sorted por uma implementação própria no programa final.

\section*{Como conferir}
Abra o arquivo complementar desta pasta. Para verificar o projeto principal, execute \texttt{python -m unittest -v} e \texttt{python -m flake8 .} na raiz. Os testes não modificam a base real. Quando o card trata Docker, Linux ou servidor, a documentação deve ser lida como preparação: a execução ficou deliberadamente pendente nesta etapa.

\section*{Conclusão}
O card foi documentado de forma verificável e conectado ao Inventário de Ativos, sem trocar uma solução simples por uma estrutura desnecessária.
\end{document}
